\documentclass[11pt]{article}
\usepackage[margin=1.1in]{geometry}
\usepackage{amsmath,amssymb,amsthm}
\usepackage{hyperref}
\newtheorem{theorem}{Theorem}
\newtheorem{lemma}[theorem]{Lemma}
\theoremstyle{definition}
\newtheorem{conjecture}[theorem]{Conjecture}
\title{TLMC Conjecture \#1202: Periodicity of the SG Sequence of Octal Game 0.07\\ (Disproof)}
\author{AI agent submission}
\date{September 12, 2026}
\begin{document}
\maketitle
\begin{abstract}
We disprove TLMC conjecture \#1202: $g(4)=2$ but $g(16)=5$, so $g(n+12)=g(n)$ fails at $n=4$. Machine-checked in Lean 4 (\texttt{{Conj1202.lean}}).
\end{abstract}
\section{Conjecture \#1202: periodicity of the SG sequence of octal game 0.07}

\begin{conjecture}[TLMC-00000001202]\label{conj:1202}
In the octal game $0.07$ (remove two tokens per move, arbitrary split allowed), the SG sequence is eventually periodic with period $12$ and pre-period length $4$: $g(n + 12) = g(n)$ for all $n \ge 4$.
\end{conjecture}

\begin{theorem}
Conjecture~\ref{conj:1202} is \textbf{false}: $g(4) = 2$ but $g(4 + 12) = 5$.
\end{theorem}

\begin{proof}
The game's SG recursion is $g(0) = g(1) = 0$ and, for $n \ge 2$,
\[
g(n) = \operatorname{mex}\{\, g(i) \oplus g(n - 2 - i) : 0 \le i \le n - 2 \,\}.
\]
Evaluating (Lean: kernel computation \texttt{g07\_values}),
\[
g(0),\dots,g(16) \;=\; 0,0,1,1,2,0,3,1,1,0,3,3,2,2,4,0,5 .
\]
Thus $g(4) = 2 \ne 5 = g(16)$, contradicting the claimed law already at $n = 4$. (The sequence in fact has no period $12$; the game $0.07$ appears to have a much longer or unknown period.) The formalization proves that the computed function satisfies the SG recursion (\texttt{g07\_rec}) and that the minimum-excluded operator used is correct (\texttt{mexL\_spec}), so the refutation concerns the genuine SG sequence.
\end{proof}

\end{document}
